% Decision theory: two posteriors, same mean, different uncertainty
\begin{figure}[htbp]
  \centering
  \begin{tikzpicture}[x=3cm, y=1.1cm]
    % axes
    \draw[-{Latex[length=2.5mm]}, thick] (-1.3,0) -- (2.6,0)
      node[right, font=\small] {CLV};
    \draw[-{Latex[length=2mm]}, thick] (-1.3,0) -- (-1.3,3.1)
      node[above, font=\small] {density};

    % cost threshold c
    \draw[accentorange, very thick, dashed] (0.45,0) -- (0.45,2.9);
    \node[accentorange, font=\small, anchor=south] at (0.45,2.9) {cost $c$};

    % Customer A: narrow, high-confidence (mostly above c)
    \draw[bayesblue, very thick, domain=0.35:1.65, samples=120, smooth]
      plot (\x, {2.6*exp(-((\x-1.0)^2)/(2*0.17*0.17))});
    \node[bayesblue, font=\footnotesize, anchor=west] at (1.25,2.2) {Customer A};
    \node[bayesblue, font=\footnotesize, anchor=west] at (1.25,1.8) {(high confidence)};

    % Customer B: wide, uncertain (straddles c)
    \draw[tealgreen, very thick, domain=-1.2:3.2, samples=160, smooth]
      plot (\x, {0.95*exp(-((\x-1.0)^2)/(2*0.62*0.62))});
    \node[tealgreen, font=\footnotesize, anchor=west] at (1.7,0.9) {Customer B};
    \node[tealgreen, font=\footnotesize, anchor=west] at (1.7,0.5) {(uncertain)};

    % shared mean marker
    \draw[gray, dotted] (1.0,0) -- (1.0,2.7);
    \node[gray, font=\footnotesize, anchor=north] at (1.0,-0.05) {shared mean};
  \end{tikzpicture}
  \caption[Decision theory under uncertainty]{Two customers with the \emph{same} posterior-mean CLV but different uncertainty. A point-estimate rule treats them identically. The decision-theoretic rule targets customer $i$ when $P(\text{CLV}_i > c) > \tau$: customer A places almost all posterior mass above the cost $c$ and is a confident target, whereas customer B's mass straddles $c$, so the probabilistic rule appropriately discounts the apparent value.}
  \label{fig:decision}
\end{figure}
